\documentclass[12pt]{article}
\usepackage[utf8]{inputenc}
\usepackage[T1]{fontenc}
\usepackage{hyperref}
\title{Margany Language Resources for K-12 Curriculum Development}
\author{Compiled by Claude Sonnet 5.5 (claude-sonnet-5-5) for Andr\'as Kornai}
\date{}
\begin{document}
\maketitle

\section{Introduction}
Margany (ISO 639-3 zmc; Glottocode marg1253) is a Pama-Nyungan language of south-west Queensland, Australia. It was spoken from around Quilpie south to Thargomindah and is also recorded as Mardigan, Marukanji and Marganj. It is closely related to Gunya (Gunja) of the Cunnamulla region, and the two are treated together in the main published description. ASJP lists it as extinct since 1990. The searches found no school programme or revitalization curriculum for it.

\section{Primary Educational Resources}
No primers, textbooks or school programmes for Margany turned up in the searches. A revitalization curriculum would have to be built from archival and published linguistic documentation, in consultation with descendant communities.

\section{Linguistic and Reference Materials}
\begin{itemize}
\item Gavan Breen, ``Margany and Gunya'', in Dixon and Blake (eds.), Handbook of Australian Languages, vol.\ 2 (1981). This is the main published grammar and wordlist. Bibliographic records only: \url{https://phoible.org/sources/breen_margany_1981}, \url{https://phoible.org/sources/132305}, \url{https://wals.info/refdb/record/Breen-1981}
\item AIATSIS finding aid for the Breen papers (MS 113). It lists Margany and Gunya word lists, a classified vocabulary, a draft grammar and tape transcriptions: \url{https://aiatsis.gov.au/sites/default/files/catalogue_resources/MS113.PDF}
\item AIATSIS AUSTLANG entry for the language: \url{https://collection.aiatsis.gov.au/austlang/language/d43}
\item Grambank entry: \url{https://grambank.clld.org/languages/marg1253}
\item ASJP wordlist entry: \url{https://asjp.clld.org/languages/MARGANY_2}
\end{itemize}

\section{Religious and Community Resources}
The searches found no scripture translations or community organisations specific to Margany.

\section{Digital and Audio Resources}
\begin{itemize}
\item State Library of Queensland, Word Week blog for the International Year of Indigenous Languages 2019 (week thirty-eight). I did not confirm whether it features Margany words: \url{https://www.slq.qld.gov.au/blog/2019-international-year-indigenous-languages-word-week-week-thirty-eight}
\item Handbook of the Cosmos (hgcosmos.org) Margany page: \url{https://www.hgcosmos.org/assets/culture/AZ/M/Margany.pdf}
\end{itemize}
The archived tape recordings listed in the AIATSIS finding aid (Cunnamulla, Eulo and Quilpie, 1967--1976) are the likely source of audio. I did not confirm that any are available online.

\section{Curriculum Development Context}
The language has no living first-language speakers. Revitalization would rely on Breen's documentation and the AIATSIS holdings. Access to and use of those materials, and any teaching, should be agreed with the Aboriginal communities of the Quilpie, Eulo and Cunnamulla region.

\end{document}